% SPDX-FileCopyrightText: 2026 Claire Tam <claire.tam@student.adelaide.edu.au>
% SPDX-FileCopyrightText: 2026 fractuscontext <106440141+fractuscontext@users.noreply.github.com>
%
% SPDX-License-Identifier: LPPL-1.3c

\documentclass[a4paper,11pt]{article}
\usepackage[margin=1in]{geometry}
\usepackage{ifxetex,ifluatex}
\usepackage{etoolbox}
\usepackage[T1]{fontenc}
\usepackage[utf8]{inputenc}
\usepackage{helvet}
\renewcommand{\familydefault}{\sfdefault}

\usepackage{user-persona}

\begin{document}

\begin{persona}{Tyler Durden}
  % Placeholder for photo
  \personaphoto{tyler-placeholder.jpg} 
  
  \personaquote{``If the university portal crashes while I am on hold 
    with Centrelink and trying to sort out Mum's NDIS plan, my entire 
    week's schedule completely collapses.''} % chktex 38
  
  \personafact{Age}{26}
  \personafact{Gender}{Non-binary (They/Them)}
  \personafact{Location}{Melbourne, Australia}
  \personafact{Education}{Part-time B.Sc (Honours), High Distinction}
  \personafact{Status}{NDIS Participant, Centrelink, Carer}
  \personafact{Work}{PT Retail \& Advocacy Spokesperson}
  \personafact{Neurotype}{Autistic / ADHD}
  \personafact{Energy}{Exhausted (Low/Fluctuating)}
  
  \personasection{Goals \& Motivations}{PersonaGreen}{
    \item Maintain a High Distinction academic average despite extremely limited energy slots.
    \item Successfully advocate for neurodivergent and carer communities at the university level.
    \item Manage complex caregiving duties for their mother without sacrificing their Honours research.
    \item Secure a PhD scholarship in Disability Studies to influence policy from within the system.
  }
  
  \personasection{Behaviours}{PersonaBlue}{
    \item Relies on strict, militant ``time-blocking'' to juggle uni, work, and caregiving.
    \item Prefers asynchronous, text-based communication over unpredictable live phone calls due to sensory load.
    \item Completes heavyweight academic drafting in hospital waiting rooms or during brief quiet periods at night.
    \item Uses ``body doubling'' apps and focus timers to stay on task during brain-fog episodes.
    \item Audits every new software tool for keyboard accessibility and dark-mode support before adoption.
  }
  
  \personasection[right]{Pain Points}{PersonaRed}{
    \item Bureaucratic red tape: Navigating rigid university extension policies vs.\ unpredictable NDIS/carer crises.
    \item Cognitive fatigue from constantly having to self-advocate for basic accessibility accommodations.
    \item Severe burnout when group-mates fail to respect highly structured, pre-planned meeting times.
    \item Financial anxiety due to the unreliable nature of government support payments and medical costs.
  }
  
  \personacard[right]{Tools \& Tech}{PersonaDark}{
    \personapill{Notion}
    \personapill{Google Calendar}
    \personapill{Slack}
    \personapill{Zotero}
    \personapill{LiquidText}
    \personabrand[fa]{\faApple}
    \personabrand[fa]{\faAccessibleIcon}
    \personabrand[fa]{\faCoffee}
    \personabrand[fa]{\faBrain}
  }

  % Context cards on the left
  \personacontext{Frequency}{Scheduled Bursts (highly dependent on daily capacity)}
  \personacontext{Motivation}{Academic survival and pushing systemic change}
  \personacontext{Setting}{Home office, hospital waiting rooms, or commuting on public transport}
  
  \personacard[left]{Life Events}{PersonaGrey}{
    \textbf{2022:} Diagnosed with burnout.\\
    \textbf{2023:} Became full-time carer.\\
    \textbf{2024:} Started Honours project.
  }

  \personasection[span]{Core Values \& Advocacy}{PersonaGreen}{
    % not everyone uses aussie spelling, but this is a persona for an Australian student so it makes sense to use it here
    \item \textbf{Intersectionality:} Ensuring advocacy includes the most marginalised voices in the community.
    \item \textbf{Transparency:} Pushing for open, clear, and honest communication from institutional portals.
    \item \textbf{Autonomy:} Designing systems that empower users to manage their own needs without friction.
  }

  \personasection[span]{Accessibility Requirements}{PersonaGrey}{
    \item High contrast text for reading in brightly lit hospital environments.
    \item Logical heading structures for screen reader compatibility during high-fatigue periods.
    \item Predictable interaction patterns to reduce cognitive load during multi-tasking.
  }
  
\end{persona}

\end{document}
